\documentclass[11pt]{article}
\usepackage{mathblatt}
\usepackage{multicol}


\begin{document}
\pfheftstil
\pfheftkopf{Koerper \textperiodcentered{} Lösungen}{}
\begin{pfloesung}{}
\lz{1.}{6 cm | 36 cm²; 3,2 cm | 10,24 cm²; 29 cm | 841 cm²}{}{}
\end{pfloesung}
\begin{pfloesung}{}
\lz{2.}{2 Dreiecke und 3 Rechtecke; 4 Dreiecke; 2 Kreise und 1 Rechteck}{}{}
\end{pfloesung}
\begin{pfloesung}{}
\lz{3.}{Durchmesser}{}{}
\end{pfloesung}
\begin{pfloesung}{}
\lz{4.}{Radius $4$ cm}{Durchmesser}{}
\end{pfloesung}
\begin{pfloesung}{}
\lz{5.}{Volumen}{}{}
\end{pfloesung}
\begin{pfloesung}{}
\lz{6.}{$50$ cm und $70$ cm}{}{}
\end{pfloesung}
\begin{pfloesung}{}
\lz{7.}{Mantellinie}{}{}
\end{pfloesung}
\begin{pfloesung}{}
\lz{8.}{auch wenn das Prisma liegt}{Grundfläche: Grundfläche: eines der beiden gleichen Dreiecke; die Deckfläche ist das andere}{}
\end{pfloesung}
\begin{pfloesung}{}
\lz{9.}{Seitenhöhe}{}{}
\end{pfloesung}
\begin{pfloesung}{}
\lz{10.}{Höhe von der Spitze senkrecht zur Mitte der Grundfläche, von dort zur Mitte einer Grundkante , zurück zur Spitze}{halbe Grundkante, 3 m}{}
\end{pfloesung}
\begin{pfloesung}{}
\lz{11.}{$36$ cm³}{$6 \cdot 3 \cdot 2$ $\Rightarrow$ 36}{}
\end{pfloesung}
\begin{pfloesung}{}
\lz{12.}{V = $\tfrac{1}{3}$ · 25 cm² · 6 cm = 50 cm³}{}{}
\end{pfloesung}
\begin{pfloesung}{}
\lz{13.}{27 cm³}{(3 cm)³ = 27 cm³}{}
\end{pfloesung}
\begin{pfloesung}{}
\lz{14.}{(3): die Deckfläche des Zylinders wird mitgezählt}{}{}
\end{pfloesung}
\begin{pfloesung}{}
\lz{15.}{72$\pi$ (V = $\pi$ · 3² · 8)}{}{}
\end{pfloesung}
\begin{pfloesung}{}
\lz{16.}{V = $\tfrac{1}{3}$ · 4² · 6 cm³ = 32 cm³}{}{}
\end{pfloesung}
\begin{pfloesung}{}
\lz{17.}{3 000 Würfel (20 dm · 10 dm · 15 dm = 3 000 dm³)}{}{}
\end{pfloesung}
\begin{pfloesung}{}
\lz{18.}{$105$ cm³}{Grundfläche: $7 \cdot 3$ $\Rightarrow$ 21 cm²; $21 \cdot 5$ $\Rightarrow$ 105}{}
\end{pfloesung}
\begin{pfloesung}{}
\lz{19.}{$63\,000$ cm³}{Grundfläche: $90 \cdot 20 : 2$ $\Rightarrow$ 900 cm²; $900 \cdot 70$ $\Rightarrow$ 63000}{}
\end{pfloesung}
\begin{pfloesung}{}
\lz{20.}{V = 6,5 · 4,5 · 0,3 $\approx$ 8,8 m³}{}{}
\end{pfloesung}
\begin{pfloesung}{}
\lz{21.}{$\tfrac{9 \pi}{2} \approx 14{,}1$ l}{$\tfrac{4}{3} \cdot \pi \cdot 15^3$ $\Rightarrow$ $\approx$ 14\,137{,}2 cm³}{}
\end{pfloesung}
\begin{pfloesung}{}
\lz{22.}{das sind etwa $350$ ml}{$\pi \cdot 3{,}5^2 \cdot 9$ $\Rightarrow$ $\approx$ 346{,}4 cm³ $\approx$ 346 ml}{}
\end{pfloesung}
\begin{pfloesung}{}
\lz{23.}{$\approx 5{,}3$ l}{$\pi \cdot 11^2 \cdot 14$ $\Rightarrow$ $\approx$ 5\,321{,}9 cm³}{}
\end{pfloesung}
\begin{pfloesung}{}
\lz{24.}{$\approx 415$ ml}{$\tfrac{1}{3} \cdot \pi \cdot 6^2 \cdot 11$ $\Rightarrow$ $\approx$ 414{,}7 cm³}{}
\end{pfloesung}
\begin{pfloesung}{}
\lz{25.}{862 125 cm³ ($\approx$ 0,86 m³)}{5 225 cm² · 165 cm}{}
\end{pfloesung}
\begin{pfloesung}{}
\lz{26.}{288$\pi$ $\approx$ 904,8 cm³}{6 cm; $\tfrac{4}{3}$ · $\pi$ · 6³}{}
\end{pfloesung}
\begin{pfloesung}{}
\lz{27.}{wahr}{V = $\pi$ · 3,2² · 7 = 71,68$\pi$ $\approx$ 225,2 cm³ > 200 ml; Grundfläche: 10,24$\pi$ $\approx$ 32,2 cm²; 1 cm³ = 1 ml}{}
\end{pfloesung}
\begin{pfloesung}{}
\lz{28.}{wahr}{V = $\pi$ · 4² · 8 = 128$\pi$ $\approx$ 402,1 cm³ $\approx$ 400 ml; Grundfläche: 16$\pi$ $\approx$ 50,3 cm²; 1 cm³ = 1 ml}{}
\end{pfloesung}
\begin{pfloesung}{}
\lz{29.}{V = $\pi$ · 29² · 95 $\approx$ 250 998 cm³ $\approx$ 251 l}{1 l = 1 dm³ = 1000 cm³}{}
\end{pfloesung}
\begin{pfloesung}{}
\lz{30.}{Prisma}{}{}
\end{pfloesung}
\begin{pfloesung}{}
\lz{31.}{$8$ cm}{Kreis: 3 cm; Kreisausschnitt: Kreisausschnitt: Radius; die Mantellinie}{}
\end{pfloesung}
\begin{pfloesung}{}
\lz{32.}{Quader}{}{}
\end{pfloesung}
\begin{pfloesung}{}
\lz{33.}{ja}{}{}
\end{pfloesung}
\begin{pfloesung}{}
\lz{34.}{zwischen B und D liegt genau eine Fläche}{}{}
\end{pfloesung}
\begin{pfloesung}{}
\lz{35.}{Prisma}{}{}
\end{pfloesung}
\begin{pfloesung}{}
\lz{36.}{Pyramide}{}{}
\end{pfloesung}
\begin{pfloesung}{}
\lz{37.}{mittleres Quadrat der senkrechten Dreierreihe}{}{}
\end{pfloesung}
\begin{pfloesung}{}
\lz{38.}{$4 \pi \approx 12{,}6$ cm, also $5$ cm × $12{,}6$ cm}{Umfang: $2 \cdot \pi \cdot 2$}{}
\end{pfloesung}
\begin{pfloesung}{}
\lz{39.}{drittes Netz; Länge 6,4$\pi$ $\approx$ 20,1 cm, Breite 7 cm}{Umfang 2 · $\pi$ · 3,2}{}
\end{pfloesung}
\begin{pfloesung}{}
\lz{40.}{$6$ cm}{$90 : (5 \cdot 3)$ $\Rightarrow$ 6}{}
\end{pfloesung}
\begin{pfloesung}{}
\lz{41.}{$3\,000 : (\pi \cdot 10^2) \approx 9{,}5$ cm}{10 cm, V = 3\,000 cm³; h}{}
\end{pfloesung}
\begin{pfloesung}{}
\lz{42.}{$\approx 6{,}5$ cm}{$\sqrt{2\,000 : (\pi \cdot 15)}$; erst teilen, dann die Wurzel}{}
\end{pfloesung}
\begin{pfloesung}{}
\lz{43.}{h = 125/(2$\pi$) $\approx$ 19,9 cm}{1 l = 1000 cm³; 1000 : (16$\pi$); Grundfläche: 16$\pi$ $\approx$ 50,3 cm²}{}
\end{pfloesung}
\begin{pfloesung}{}
\lz{44.}{r = √(425/(7$\pi$)) $\approx$ 4,40 cm}{425 : ($\pi$ · 7) $\Rightarrow$ 425/(7$\pi$) $\approx$ 19,33}{}
\end{pfloesung}
\begin{pfloesung}{}
\lz{45.}{Quader und Dreiecksprisma}{Wohnteil; Dach}{}
\end{pfloesung}
\begin{pfloesung}{}
\lz{46.}{$96$ dm³}{$8 \cdot 4 \cdot 2 + 8 \cdot 2 \cdot 2$ $\Rightarrow$ 64 + 32}{}
\end{pfloesung}
\begin{pfloesung}{}
\lz{47.}{130 cm³ (150 $-$ 2 · 10)}{}{}
\end{pfloesung}
\begin{pfloesung}{}
\lz{48.}{$\approx 1{,}76$ kg}{1{,}1 cm, innen r = 1{,}0 cm; $\pi \cdot (1{,}1^2 - 1^2) \cdot 300$ $\Rightarrow$ $\approx$ 197{,}9 cm³; 1\,761 g}{}
\end{pfloesung}
\begin{pfloesung}{}
\lz{49.}{$a : 2$, die Kugel hat $\tfrac{\pi}{6} \cdot a^3$}{6 cm; 904{,}8 cm³. b) Würfel 1\,728 cm³; 823{,}2 cm³ $\approx$ 47{,}6 \%. c) Mit Kante a ist r}{}
\end{pfloesung}
\begin{pfloesung}{}
\lz{50.}{r\_innen = 13,5 m $-$ 0,03 m = 13,47 m; V = $\tfrac{4}{3}$ · $\pi$ · 13,47³ m³ $\approx$ 10 237,44 m³}{}{}
\end{pfloesung}
\begin{pfloesung}{}
\lz{51.}{V = 7,1 · 180 = 1 278 cm³ je Rohr; 130 · 1 278 · 0,7 g $\approx$ 116 kg < 250 kg: reicht}{}{}
\end{pfloesung}
\begin{pfloesung}{}
\lz{52.}{$2\,142{,}9$ cm³}{Karton: $24 \cdot 24 \cdot 9$ $\Rightarrow$ 5\,184 cm³; Torte: $\pi \cdot 11^2 \cdot 8$ $\Rightarrow$ $\approx$ 3\,041{,}1 cm³; Unterschied: 5\,184 - 3\,041{,}1 $\Rightarrow$ 2142{,}9}{}
\end{pfloesung}
\begin{pfloesung}{}
\lz{53.}{51 350,25$\pi$ $\approx$ 161 322 cm³}{$\pi$ · 30,5² · 81 $\Rightarrow$ 75 350,25$\pi$ $\approx$ 236 719,8 cm³; $\tfrac{1}{3}$ · $\pi$ · 30² · 80 $\Rightarrow$ 24 000$\pi$ $\approx$ 75 398,2 cm³}{}
\end{pfloesung}
\begin{pfloesung}{}
\lz{54.}{M = 4 · ½ · 7 · 8,73 cm² = 122,22 cm²}{}{}
\end{pfloesung}
\begin{pfloesung}{}
\lz{55.}{O = 2 · (275 + 225 + 99) cm² = 1 198 cm²}{}{}
\end{pfloesung}
\begin{pfloesung}{}
\lz{56.}{$72 \pi \approx 226{,}2$ cm²}{$2 \cdot \pi \cdot 4 \cdot 9$}{}
\end{pfloesung}
\begin{pfloesung}{}
\lz{57.}{$\pi \cdot 4 \cdot 11 \approx 138{,}2$ cm²}{M}{}
\end{pfloesung}
\begin{pfloesung}{}
\lz{58.}{138,88$\pi$ $\approx$ 436,3 m²}{$\pi$ · 6,4 · 21,7 $\Rightarrow$ 138,88$\pi$ m²}{}
\end{pfloesung}
\begin{pfloesung}{}
\lz{59.}{1~Mantelfläche mit Kosten, 2~Volumen direkt, 3~Mantelfläche mit Kosten, 4~Volumen direkt, 5~Mantelfläche mit Kosten, 6~Volumen direkt}{}{}
\end{pfloesung}
\begin{pfloesung}{}
\lz{60.}{V ~ r³, 2³ = 8}{}{}
\end{pfloesung}
\begin{pfloesung}{}
\lz{61.}{Nein}{2a · 2b · 2c = 8 · abc – das Volumen wird achtmal so groß}{}
\end{pfloesung}
\begin{pfloesung}{}
\lz{62.}{V ~ r²: 3² = 9 Stück}{}{}
\end{pfloesung}
\begin{pfloesung}{}
\lz{63.}{vervierfacht (V ~ r²)}{}{}
\end{pfloesung}
\begin{pfloesung}{}
\lz{64.}{Nein}{das Volumen verdoppelt sich (V = a · b · 2c).}{}
\end{pfloesung}
\begin{pfloesung}{}
\lz{65.}{Der Radius steht im Quadrat: Der doppelte Radius gibt eine viermal so große Grundfläche, die Höhe bleibt gleich}{}{}
\end{pfloesung}
\begin{pfloesung}{}
\lz{66.}{Das Kegelvolumen ist ein Drittel von Grundfläche mal Höhe, das Zylindervolumen Grundfläche mal Höhe}{}{}
\end{pfloesung}
\begin{pfloesung}{}
\lz{67.}{das Achtfache, weil $2^3 = 8$}{A: 33{,}5 cm³, B: V $\approx$ 268{,}1 cm³; 268{,}1 : 33{,}5 $\Rightarrow$ $\approx$ 8}{}
\end{pfloesung}
\begin{pfloesung}{}
\lz{68.}{nein}{Volumen vervierfacht sich (319,58$\pi$ l $\approx$ 1004 l); (2r)² = 4r²}{}
\end{pfloesung}
\begin{pfloesung}{}
\lz{69.}{sechs Rechtecke, je zwei gleich: $4 \times 2$, $4 \times 1$ und $2 \times 1$ , so aneinander, dass sie zu einem Quader falten}{in cm}{}
\end{pfloesung}
\begin{pfloesung}{}
\lz{70.}{zwei Dreiecke mit den Seiten $8$ cm, $5$ cm, $5$ cm}{}{}
\end{pfloesung}
\begin{pfloesung}{}
\lz{71.}{Rückfläche 15 cm × 12 cm rechts an die Grundfläche; zweites Trapez unten an die Grundfläche}{}{}
\end{pfloesung}
\begin{pfloesung}{}
\lz{72.}{alle vier Seiten des Quadrats $4$ cm}{}{}
\end{pfloesung}
\begin{pfloesung}{}
\lz{73.}{O = 2 · (15 · 20 + 10 · 20 + 10 · 15) = 1 300 cm²; Netz mit Rechtecken 4 × 3, 4 × 2, 3 × 2 cm}{}{}
\end{pfloesung}
\begin{pfloesung}{}
\lz{74.}{Netz mit drei Rechtecken und zwei gleichseitigen Dreiecken; a = 1,50 m}{}{}
\end{pfloesung}
\begin{pfloesung}{}
\lz{75.}{Netz um Rechtecke 57 cm × 165 cm und 110 cm × 165 cm ergänzt; a = 110 cm, b = 57 cm, d = 165 cm}{}{}
\end{pfloesung}
\begin{pfloesung}{}
\lz{76.}{$20 \pi \approx 62{,}8$ cm³}{$\pi \cdot 2^2 \cdot 5$}{}
\end{pfloesung}
\begin{pfloesung}{}
\lz{77.}{Rechteck, a = 8$\pi$ $\approx$ 25,1 cm, b = 8 cm}{2 · $\pi$ · r (Umfang); b = h}{}
\end{pfloesung}
\begin{pfloesung}{}
\lz{78.}{Tiefe $4$ cm schräg als $2$ cm}{Vorderseite: $6$ cm $\times 3$ cm in wahrer Größe; Kästchendiagonalen}{}
\end{pfloesung}
\begin{pfloesung}{}
\lz{79.}{Spitze in der Mitte des Deckels}{Grundkreis auf dem Boden, er berührt alle vier Seitenwände , im Schrägbild als Ellipse; Durchmesser: 50 cm; Schnitt der Diagonalen}{}
\end{pfloesung}
\begin{pfloesung}{}
\lz{80.}{Quadrat mit $5$ cm Seitenlänge, an jeder Seite ein gleichschenkliges Dreieck mit Grundseite $5$ cm und Höhe $7$ cm}{}{}
\end{pfloesung}
\begin{pfloesung}{}
\lz{81.}{oben der Kreis als Ellipse mit dem Radius $2{,}5$ cm}{Kegel mit der Spitze unten}{}
\end{pfloesung}
\begin{pfloesung}{}
\lz{82.}{Skizze: Ellipse auf der Bodenfläche, Spitze mittig oben; 61 cm × 61 cm × 81 cm; 31 cm schmaler als Durchmesser 60 cm, 101 cm Höhe unnötig groß}{Durchmesser: 2 · 30 $\Rightarrow$ 60 cm; Höhe: 80 cm}{}
\end{pfloesung}
\begin{pfloesung}{}
\lz{83.}{flache Ellipse durch den Mittelpunkt als Äquator}{Kreis mit dem Radius: 2 cm; hintere Hälfte gestrichelt}{}
\end{pfloesung}
\begin{pfloesung}{}
\lz{84.}{Höhe von der Spitze zum Diagonalenschnittpunkt, h = 0,68 m; Grundkante 0,80 m}{}{}
\end{pfloesung}
\begin{pfloesung}{\textbf{Prüfstein}}
\lz{a)}{552,25$\pi$ $\approx$ 1734,9 m³}{$\pi$ · 4,7² · 25 $\Rightarrow$ 552,25$\pi$ m³; 22,09 m²}{}
\lz{b)}{$\approx$ 2539,66 € (rund 2540 €)}{40,42$\pi$ $\approx$ 126,98 m²}{}
\lz{c)}{25 + √51,87 $\approx$ 32,2 m}{√(8,6² $-$ 4,7²) $\Rightarrow$ √51,87 $\approx$ 7,20 m}{}
\lz{d)}{falsch}{V\_Kegel = $\tfrac{1}{3}$ · $\pi$ · (3r)² · h = 3 · $\pi$ · r² · h = 3 · V\_Zylinder $\Rightarrow$ dreifaches Volumen; (3r)² = 9r²; 9 : 3 $\Rightarrow$ 3}{}
\end{pfloesung}
\end{document}